% The Appendix section (PUB-REF-0001 CP-16: codes from spec/iword/
% reveal_codes.h, the file format from spec/iword/wp5_format.h, messages,
% keyboards; model: WordPerfect 5.1 Reference Appendix A..T, each on a
% recto, docs/pubs/originals/wordperfect-51.md s3).  Every code, message and
% key alias below is registered (\WPCode, \WPMessage, \WPAlias) and graded
% both ways against the surface measured from WORD.EXE by `make
% test-pubs-iword` (harness/pubs/iword_probe.c): the codes from the locked
% table's text, the messages by driving the program into each one, the keys
% by pressing each.
% GATES: test-pubs-iword test-iword-wp5-roundtrip test-iword-reveal-codes test-iword-face
\Gates{test-pubs-iword test-iword-wp5-roundtrip test-iword-reveal-codes test-iword-face}
\WPHead{Appendix Contents}
\begin{tabular}{@{}ll@{}}
Appendix A & Characters\\
Appendix B & Codes\\
Appendix C & Messages and Prompts\\
Appendix D & Document File Format\\
Appendix E & Initial Settings\\
Appendix F & Keyboard\\
Appendix G & Menus\\
Appendix H & Limits\\
\end{tabular}

%% ---------------------------------------------------------------------------
\WPAppendix{A}{Characters}
InitechWord takes the characters of the keyboard that print: the space and
the 94 characters from \Scr{!} to \Scr{\textasciitilde}.  Each is stored in
the document as the one byte of its ASCII code, as WordPerfect 5.1 stores
it.  The table gives each character with its decimal code.

\par\addvspace{6pt}
\directlua{
  local out = {}
  for c = 32, 126 do
    local ch = string.char(c)
    local shown
    if c == 32 then shown = "sp"
    elseif ch:match("[A-Za-z0-9]") then shown = ch
    else shown = "\string\\char" .. c .. "\string\\relax " end
    out[\string#out + 1] = c .. " & \string\\texttt{" .. shown .. "}"
  end
  local rows = {}
  local per = 19
  for r = 1, per do
    local cells = {}
    for col = 0, 4 do
      local i = r + col * per
      if out[i] then cells[\string#cells + 1] = out[i] else cells[\string#cells + 1] = " & " end
    end
    rows[\string#rows + 1] = table.concat(cells, " & ") .. "\string\\\string\\"
  end
  tex.sprint("\string\\noindent{\string\\small\string\\begin{tabular}{@{}rl@{\string\\quad}rl@{\string\\quad}rl@{\string\\quad}rl@{\string\\quad}rl@{}}" .. table.concat(rows) .. "\string\\end{tabular}}")
}
\par\addvspace{6pt}

The keys that are not characters -- the function keys, Enter, Tab and the
others -- do not put characters in the document; they put codes there or
move the cursor (see Appendix F).

\Notes
\Topic{Characters From Other Documents}
A WordPerfect 5.1 document can hold the hard space, shown in Reveal Codes
as \Code{[ ]}: a space at which a line never wraps.  InitechWord keeps hard
spaces in documents it retrieves; it has no key that types one.

%% ---------------------------------------------------------------------------
\WPAppendix{B}{Codes}
These are the codes Reveal Codes shows, in the order of the alphabet of
their names.  A pair of codes is listed under its on code (capitals) and its
off code (lower case).  Codes marked \textit{kept} are ones InitechWord keeps
in a document it retrieves but does not put in documents itself (see
\InRef{Codes}).
\par\addvspace{8pt}
\WPCode{[ ]}{Hard space: a space at which a line does not wrap (kept).}
\WPCode{[?]}{A code of a WordPerfect 5.1 document that Reveal Codes has no
name for (kept).}
\WPCode{[->Indent]}{Indent, F4.  Shown with an arrow, \Code{[\rarr Indent]}.}
\WPCode{[Block]}{Not a code: shows where Block was turned on, while it is on.}
\WPCode{[BOLD]}{Bold on, F6.}
\WPCode{[bold]}{Bold off.}
\WPCode{[Center]}{Center a line, from a WordPerfect document (kept).}
\WPCode{[DBL UND]}{Double underline on, Ctrl-F8,2,3.}
\WPCode{[dbl und]}{Double underline off.}
\WPCode{[DSRt]}{Deletable soft return, put in by Justification in the middle
of a line.}
\WPCode{[EXT LARGE]}{Extra large on, Ctrl-F8,1,7.}
\WPCode{[ext large]}{Extra large off.}
\WPCode{[FINE]}{Fine on, Ctrl-F8,1,3.}
\WPCode{[fine]}{Fine off.}
\WPCode{[Flsh Rgt]}{Flush right, from a WordPerfect document (kept).}
\WPCode{[HPg]}{Hard page, Ctrl-Enter.}
\WPCode{[HRt]}{Hard return, Enter.}
\WPCode{[ITALC]}{Italic on, Ctrl-F8,2,4.}
\WPCode{[italc]}{Italic off.}
\WPCode{[Just:Center]}{Justification center, Shift-F8,1,3,2.}
\WPCode{[Just:Full]}{Justification full, Shift-F8,1,3,4.}
\WPCode{[Just:Left]}{Justification left, Shift-F8,1,3,1.}
\WPCode{[Just:Right]}{Justification right, Shift-F8,1,3,3.}
\WPCode{[L/R Mar:...]}{Left and right margins, Shift-F8,1,7, with the two
margins in inches: \Code{[L/R Mar:1.5",2"]}.}
\WPCode{[LARGE]}{Large on, Ctrl-F8,1,5.}
\WPCode{[large]}{Large off.}
\WPCode{[Ln Spacing:...]}{Line spacing, Shift-F8,1,6, with the spacing:
\Code{[Ln Spacing:2]}.}
\WPCode{[OUTLN]}{Outline on, Ctrl-F8,2,5.}
\WPCode{[outln]}{Outline off.}
\WPCode{[REDLN]}{Redline on, Ctrl-F8,2,8.}
\WPCode{[redln]}{Redline off.}
\WPCode{[SHADW]}{Shadow on, Ctrl-F8,2,6.}
\WPCode{[shadw]}{Shadow off.}
\WPCode{[SM CAP]}{Small capitals on, Ctrl-F8,2,7.}
\WPCode{[sm cap]}{Small capitals off.}
\WPCode{[SMALL]}{Small on, Ctrl-F8,1,4.}
\WPCode{[small]}{Small off.}
\WPCode{[SRt]}{Soft return: the end of a line made by word wrap.}
\WPCode{[STKOUT]}{Strikeout on, Ctrl-F8,2,9.}
\WPCode{[stkout]}{Strikeout off.}
\WPCode{[SUBSCPT]}{Subscript on, Ctrl-F8,1,2.}
\WPCode{[subscpt]}{Subscript off.}
\WPCode{[SUPRSCPT]}{Superscript on, Ctrl-F8,1,1.}
\WPCode{[suprscpt]}{Superscript off.}
\WPCode{[Tab]}{Tab, the Tab key.}
\WPCode{[UND]}{Underline on, F8.}
\WPCode{[und]}{Underline off.}
\WPCode{[VRY LARGE]}{Very large on, Ctrl-F8,1,6.}
\WPCode{[vry large]}{Very large off.}

\Notes
\Topic{Codes Reveal Codes Does Not Show}
WordPerfect 5.1 keeps some records in a document that Reveal Codes never
shows: one at the end of an indented paragraph, one before each hard page,
one at the start of each centered or right-justified line, and one at the
end of a centered or flush-right line.  InitechWord writes them as
WordPerfect does and keeps them up to date as you edit; you never see or
delete them.

%% ---------------------------------------------------------------------------
\WPAppendix{C}{Messages and Prompts}
These are the messages, questions and prompts InitechWord shows on the
status line, in alphabetical order.  Each is shown exactly as it appears on
the screen; \textit{filename} and \textit{code} stand for the name or code
shown.

\WPMessage{(Text was not modified)}{Shown at the right end of the status
line with \Scr{Save Document? Yes (No)} when the document has not changed
since it was last saved or retrieved.}{Type \Typed{n}: there is nothing new
to save.}
\WPMessage{* Not found *}{Search or Replace found nothing more in the
direction it was going (see \InRef{Search} and \InRef{Replace}).  The cursor
is where it was before the search, or after the last replacement.}{Press any
key.  Check the search string, or search in the other direction.}
\WPMessage{-> Srch:}{The prompt of Search forward (F2), with the search
string (see \InRef{Search}); also the search string of Replace.}{Type the
text or codes to find and press F2.}
\WPMessage{<- Srch:}{The prompt of Search backward (Shift-F2), with the
search string; also the search string of Replace after Up.}{Type the text or
codes to find and press F2.}
\WPMessage{Block on}{Block is on (see \InRef{Block}).}{Move the cursor to
stretch the block, then press the key of the feature to use on it; press
Alt-F4 or Cancel to turn Block off.}
\WPMessage{Delete Block? No (Yes)}{Delete or Backspace was pressed while
Block was on.}{Type \Typed{y} to delete the block; any other key keeps it.}
\WPMessage{Delete [code]? No (Yes)}{Backspace or Delete reached a code on
the editing screen (see \InRef{Codes}).}{Type \Typed{y} to delete the code;
any other key moves the cursor past it.}
\WPMessage{Document to be Retrieved:}{The prompt of Retrieve.}{Type the
name of the document and press Enter, or press Cancel.}
\WPMessage{Document to be Saved:}{The prompt of Save, with the document's
name if it has one.}{Press Enter to save under the name shown, or type a
name and press Enter.}
\WPMessage{ERROR: Disk Full--Press Any Key to Continue}{The data volume has
no room for the document being saved.  The document is still on the
screen.}{Press any key.  Make room on the data volume, or save the document
under the name of a document you no longer need.}
\WPMessage{ERROR: Document too large}{The document to be retrieved is larger
than InitechWord can hold, or an addition would make the document on the
screen larger than that (see Appendix H).}{Press any key.  Divide the work
into smaller documents.}
\WPMessage{ERROR: File I/O error}{The file could not be made on the data
volume, or the data volume could not be read.}{Press any key and try again
with another name.}
\WPMessage{ERROR: File Not Found}{No document of that name is on the data
volume.}{Press any key, check the name, and type it again.}
\WPMessage{ERROR: Incompatible File Format}{The file is not a WordPerfect
5.1 document InitechWord can read.}{Press any key and give the name of an
InitechWord document.}
\WPMessage{ERROR: Invalid file name}{The name is not a valid file name (see
\InRef{Filenames}).}{Press any key and type the name again.}
\WPMessage{Exit WP? No (Yes)}{The second question of Exit.}{Type \Typed{y}
to leave InitechWord; \Typed{n} or Enter clears the screen for a new
document; Cancel returns to the document.}
\WPMessage{Move cursor; press Enter to retrieve.}{Text moved or copied with
Move waits to be put in the document.}{Move the cursor where the text is to
go and press Enter.}
\WPMessage{Original Document Backup? No (Yes)}{The question of Setup's
Original Document Backup (see \InRef{Setup}).}{Type \Typed{y} or \Typed{n}.}
\WPMessage{Press Reveal Codes to restore screen}{Reveal Codes is on.}{Press
Alt-F3 to return to the full editing screen.}
\WPMessage{Replace A:\wbs filename? No (Yes)}{A document with the name you
are saving under is already on the data volume.}{Type \Typed{y} to replace
it; any other key asks for a name again.}
\WPMessage{Replace with:}{The second prompt of Replace (see
\InRef{Replace}).}{Type the new text and press F2; press F2 alone to delete
what is found.}
\WPMessage{Replace? Yes (No)}{Replace with confirmation stopped at an
occurrence, highlighted as a block.}{Type \Typed{y} or press Enter to
replace it, \Typed{n} to leave it; Cancel stops replacing.}
\WPMessage{Save Document? Yes (No)}{The first question of Exit.}{Type
\Typed{y} or press Enter to save the document, \Typed{n} not to; Cancel
returns to the document.}
\WPMessage{Selection: 0}{A Setup screen waits for a choice (see
\InRef{Setup}).}{Type the number of a choice; press Exit (F7) to go back a
screen, or Cancel to return to the document.}
\WPMessage{Undelete 1 Restore; 2 Previous Deletion: 0}{Cancel (F1) was
pressed on the editing screen; the deletion shown in reverse video at the
cursor can be restored (see \InRef{Undelete}).}{Select 1 to restore it, 2
to see the deletion before it; Cancel leaves the document as it was.}
\WPMessage{Units of Measure:}{Setup asks for the letter of a unit (see
\InRef{Setup}).}{Type \Typed{i}, \Typed{c}, \Typed{p}, \Typed{w} or
\Typed{u}.}
\WPMessage{w/Confirm? No (Yes)}{The first question of Replace.}{Type
\Typed{y} to be asked at each occurrence; \Typed{n} or Enter replaces every
occurrence without asking.}

%% ---------------------------------------------------------------------------
\WPAppendix{D}{Document File Format}
An InitechWord document is a WordPerfect 5.1 document file: InitechWord
writes the file as WordPerfect 5.1 writes it, and reads the documents
WordPerfect 5.1 saves, as long as they hold only the codes listed in
Appendix B.  This appendix describes the file for programmers.  All numbers
of more than one byte are stored with the low byte first.

\WPSide{The Header}
The first 16 bytes of the file are its header:
\par\addvspace{4pt}
\noindent\begin{tabular}{@{}llp{70mm}@{}}
\textbf{Offset} & \textbf{Bytes} & \textbf{Contents}\\
0 & 4 & the byte FFh and the letters \Scr{WPC}\\
4 & 4 & where the text of the document begins (the offset of the body)\\
8 & 1 & 1, the product\\
9 & 1 & 0Ah, a document\\
10 & 1 & 0, the major version\\
11 & 1 & 1, the minor version\\
12 & 2 & 0: the document is not locked\\
14 & 2 & 0\\
\end{tabular}
\par\addvspace{4pt}

\WPSide{The Prefix}
Between the header and the body are the settings of the document, in
packets.  A new InitechWord document is saved with the 327 bytes WordPerfect
5.1 writes for a new document (the body begins at offset 327, 147h); one
byte of them records whether the cursor was at the end of the document when
it was saved.  A document that was retrieved is saved with the prefix it
came with.

\WPSide{The Body}
The body is the text and its codes, byte by byte:
\par\addvspace{4pt}
\noindent\begin{tabular}{@{}lp{76mm}@{}}
\textbf{Bytes} & \textbf{Meaning}\\
20h--7Fh & a character (Appendix A)\\
0Ah & hard return \Code{[HRt]}\\
0Ch & hard page \Code{[HPg]}\\
0Dh & soft return \Code{[SRt]}\\
83h & end of a centered or flush-right line\\
90h & deletable soft return \Code{[DSRt]}\\
A0h & hard space \Code{[ ]}\\
C0h--CFh & a fixed-length code that begins and ends with the same byte\\
D0h--FFh & a variable-length code: the byte, a sub-code, a two-byte length,
the contents, the length, the sub-code and the byte again\\
\end{tabular}
\par\addvspace{4pt}

\WPSide{Fixed-Length Codes}
\begin{tabular}{@{}llp{64mm}@{}}
\textbf{Code} & \textbf{Length} & \textbf{Meaning}\\
C1h & 9 & center, flush right or tab; the sub-code 02h is \Code{[Tab]}\\
C2h & 11 & \Code{[\rarr Indent]}\\
C3h & 3 & an attribute on: C3h, the attribute number, C3h\\
C4h & 3 & an attribute off: C4h, the attribute number, C4h\\
C6h & 6 & the end of an indented paragraph\\
\end{tabular}

\WPSide{Attribute Numbers}
\begin{tabular}{@{}rl@{\qquad}rl@{}}
0 & extra large & 8 & italic\\
1 & very large & 9 & shadow\\
2 & large & 10 & redline\\
3 & small & 11 & double underline\\
4 & fine & 12 & bold\\
5 & superscript & 13 & strikeout\\
6 & subscript & 14 & underline\\
7 & outline & 15 & small capitals\\
\end{tabular}

\WPSide{Variable-Length Codes}
\begin{tabular}{@{}llp{60mm}@{}}
\textbf{Code} & \textbf{Length} & \textbf{Meaning}\\
D0h 01h & 16 & left and right margins: the old and the new margins\\
D0h 02h & 12 & line spacing: the old and the new spacing, in 1/256 line\\
D0h 06h & 10 & justification: the old and the new, 0 left, 1 full, 2
center, 3 right\\
D4h 00h & 32 & the record WordPerfect keeps before each hard page\\
D4h 05h & 12 & the record of where a centered or right-justified line starts\\
\end{tabular}

\WPSide{Measurements}
Positions and margins are stored in WordPerfect units, 1200 to the inch,
measured from the left edge of the paper.

%% ---------------------------------------------------------------------------
\WPAppendix{E}{Initial Settings}
A new document begins with these settings.  You can change those marked
with an entry name for any part of the document.
\par\addvspace{4pt}
\noindent\begin{tabular}{@{}p{28mm}p{34mm}p{30mm}@{}}
\textbf{Setting} & \textbf{Value} & \textbf{Changed with}\\
Paper & 8.5 by 11 inches & --\\
Left and right margins & 1 inch each & Margins, Left and Right\\
Top margin & 1 inch (the first line is \Scr{Ln 1"}) & --\\
Lines on a page & 54 single-spaced lines & Page Break, Hard\\
Characters on a line & 66 & Margins, Left and Right\\
Font & Courier 10 Pitch, 10 characters to the inch & --\\
Justification & Full & Justification\\
Line spacing & 1 (single) & Line Spacing\\
Tab stops & every 0.5 inch from the left margin & --\\
Appearance and size & normal & Appearance, Size\\
\end{tabular}
\par\addvspace{4pt}
The Format: Line screen of a new document shows these values (see
\InRef{Format}).

InitechWord itself starts with these settings, kept by Setup in the file
\Scr{IWORD.SET} once you change one (see \InRef{Setup}):
\par\addvspace{4pt}
\noindent\begin{tabular}{@{}p{28mm}p{34mm}p{30mm}@{}}
\textbf{Setting} & \textbf{Value} & \textbf{Changed with}\\
Units of measure & Inches, for margins and tabs and for the status line &
Setup\\
Original document backup & No & Setup\\
Search string & empty & Search\\
\end{tabular}
\par\addvspace{4pt}

%% ---------------------------------------------------------------------------
\WPAppendix{F}{Keyboard}
This appendix lists every key InitechWord acts on.

\WPSide{The Editing Screen}
\begin{tabular}{@{}lp{60mm}@{}}
\textbf{Key} & \textbf{Action}\\
characters & typed at the cursor (Appendix A)\\
Enter & hard return\\
Ctrl-Enter & hard page\\
Tab & tab\\
Backspace & erase to the left\\
Delete & erase at the cursor\\
Up, Down, Left, Right & move the cursor\\
End & end of the line\\
Home & the Home prefix (see \InRef{Cursor Movement})\\
F1 & Cancel (in a menu or prompt); Undelete (on the editing screen)\\
Esc & Cancel (in a menu or prompt)\\
Shift-F1 & Setup\\
F2 & \rarr Search\\
Shift-F2 & \larr Search\\
Alt-F2 & Replace\\
Alt-F3 & Reveal Codes\\
F4 & \rarr Indent\\
Alt-F4 & Block\\
Ctrl-F4 & Move\\
F6 & Bold\\
F7 & Exit\\
F8 & Underline\\
Shift-F8 & Format\\
Ctrl-F8 & Font\\
F10 & Save\\
Shift-F10 & Retrieve\\
F11 & Reveal Codes, as Alt-F3\\
F12 & Block, as Alt-F4\\
\end{tabular}

\WPSide{Keys That Show a Grey Name}
These function keys name a WordPerfect feature that InitechWord does not
carry out.  Each shows the name in grey at the left of the status line and
does nothing else; the next key you press clears the name and then acts as
it always does.  The template prints the same names in grey.
\par\addvspace{4pt}
\noindent\begin{tabular}{@{}ll@{\qquad}ll@{}}
\textbf{Key} & \textbf{Name shown} & \textbf{Key} & \textbf{Name shown}\\
\WPGray{F3}{Help} & \WPGray{Alt-F1}{Thesaurus}\\
\WPGray{Alt-F5}{Mark Text} & \\
\WPGray{F9}{End Field} & \WPGray{Alt-F6}{Flush Right}\\
\WPGray{Shift-F3}{Switch} & \WPGray{Alt-F7}{Columns/Table}\\
\WPGray{Shift-F4}{Indent} & \WPGray{Alt-F8}{Style}\\
\WPGray{Shift-F5}{Date/Outline} & \WPGray{Alt-F9}{Graphics}\\
\WPGray{Shift-F6}{Center} & \WPGray{Alt-F10}{Macro}\\
\WPGray{Shift-F7}{Print} & \WPGray{Ctrl-F1}{Shell}\\
\WPGray{Shift-F9}{Merge Codes} & \WPGray{Ctrl-F2}{Spell}\\
\WPGray{Ctrl-F3}{Screen} & \WPGray{Ctrl-F5}{Text In/Out}\\
\WPGray{Ctrl-F6}{Tab Align} & \WPGray{Ctrl-F7}{Footnote}\\
\WPGray{Ctrl-F9}{Merge/Sort} & \WPGray{Ctrl-F10}{Macro Define}\\
\end{tabular}
\par\addvspace{4pt}
While Block is on, these keys do nothing (see below).

\WPSide{Keys That Do Nothing}
Ins, PgUp and PgDn do nothing, alone or with Shift, Alt or Ctrl.  Shift,
Alt and Ctrl with F11 or F12 do nothing, and neither does a function key
pressed with two or three of Shift, Alt and Ctrl.

\WPSide{Keys That Act as Other Keys}
With these keys InitechWord does not look at Shift, Alt or Ctrl: each
combination below does exactly what the key alone does.
\par\addvspace{4pt}
\noindent\begin{tabular}{@{}lll@{}}
\WPAlias{Shift-Backspace}{Backspace}, & \WPAlias{Alt-Backspace}{Backspace}, & \WPAlias{Ctrl-Backspace}{Backspace}\\
\WPAlias{Shift-Delete}{Delete}, & \WPAlias{Alt-Delete}{Delete}, & \WPAlias{Ctrl-Delete}{Delete}\\
\WPAlias{Shift-Enter}{Enter}, & \WPAlias{Alt-Enter}{Enter} & \\
\WPAlias{Shift-Tab}{Tab}, & \WPAlias{Alt-Tab}{Tab}, & \WPAlias{Ctrl-Tab}{Tab}\\
\WPAlias{Shift-Up}{Up}, & \WPAlias{Alt-Up}{Up}, & \WPAlias{Ctrl-Up}{Up}\\
\WPAlias{Shift-Down}{Down}, & \WPAlias{Alt-Down}{Down}, & \WPAlias{Ctrl-Down}{Down}\\
\WPAlias{Shift-Left}{Left}, & \WPAlias{Alt-Left}{Left}, & \WPAlias{Ctrl-Left}{Left}\\
\WPAlias{Shift-Right}{Right}, & \WPAlias{Alt-Right}{Right}, & \WPAlias{Ctrl-Right}{Right}\\
\WPAlias{Shift-End}{End}, & \WPAlias{Alt-End}{End}, & \WPAlias{Ctrl-End}{End}\\
\WPAlias{Shift-Home}{Home}, & \WPAlias{Alt-Home}{Home}, & \WPAlias{Ctrl-Home}{Home}\\
\end{tabular}
\par\addvspace{4pt}
At the \Scr{Srch:} prompt and at \Scr{Replace with:} some keys act as other
keys in the same way; they are listed under \InRef{Search} and
\InRef{Replace}.

\WPSide{While Block Is On}
The cursor keys stretch the block; \FK{Block}{Alt-F4}, \K{F12}, \K{F1} and
\K{Esc} turn it off; \K{F6}, \K{F8} and \K{Ctrl-F8} give it an attribute;
\K{Ctrl-F4} moves, copies or deletes it; \K{Delete} and \K{Backspace} ask
whether to delete it; \K{Alt-F3} and \K{F11} turn Reveal Codes on or off.
Every other key does nothing.  These also act as the key alone:
\par\addvspace{4pt}
\noindent{\small\begin{tabular}{@{}lll@{}}
\WPAlias{Block:Shift-F1}{F1}, & \WPAlias{Block:Alt-F1}{F1}, & \WPAlias{Block:Ctrl-F1}{F1}\\
\WPAlias{Block:Shift-Esc}{Esc}, & \WPAlias{Block:Alt-Esc}{Esc}, & \WPAlias{Block:Ctrl-Esc}{Esc}\\
\WPAlias{Block:Shift-Delete}{Delete}, & \WPAlias{Block:Alt-Delete}{Delete}, & \WPAlias{Block:Ctrl-Delete}{Delete}\\
\WPAlias{Block:Shift-Backspace}{Backspace}, & \WPAlias{Block:Alt-Backspace}{Backspace}, & \WPAlias{Block:Ctrl-Backspace}{Backspace}\\
\end{tabular}}
\par\addvspace{4pt}
(In this table \textit{Block:} means ``with Block on''.)

%% ---------------------------------------------------------------------------
\WPAppendix{G}{Menus}
The menus InitechWord shows, with the choices it acts on.  A choice not
listed returns to the document or leaves the screen as it is (see
\textit{Menu Options} in Getting Started).

\WPSide{Font (Ctrl-F8)}
\begin{tabular}{@{}lp{70mm}@{}}
1 Size & 1 Suprscpt, 2 Subscpt, 3 Fine, 4 Small, 5 Large, 6 Vry Large, 7 Ext Large\\
2 Appearance & 1 Bold, 2 Undln, 3 Dbl Und, 4 Italc, 5 Outln, 6 Shadw, 7 Sm Cap, 8 Redln, 9 Stkout\\
3 Normal & \\
\end{tabular}
\par With Block on: \Scr{Attribute: 1 Size; 2 Appearance: 0}, with the same
Size and Appearance menus.

\WPSide{Format (Shift-F8)}
\begin{tabular}{@{}ll@{}}
1 Line & Format: Line\\
\quad 3 Justification & 1 Left, 2 Center, 3 Right, 4 Full\\
\quad 6 Line Spacing & a spacing, then Enter\\
\quad 7 Margins & the left margin, Enter, the right margin, Enter\\
\end{tabular}

\WPSide{Move (Ctrl-F4)}
With Block off:
\par\noindent\begin{tabular}{@{}ll@{}}
1 Sentence, 2 Paragraph, 3 Page & then 1 Move, 2 Copy, 3 Delete\\
4 Retrieve & 1 Block\\
\end{tabular}
\par With Block on:
\par\noindent\begin{tabular}{@{}ll@{}}
1 Block & then 1 Move, 2 Copy, 3 Delete\\
\end{tabular}

\WPSide{Replace (Alt-F2)}
\Scr{w/Confirm? No (Yes)}, then \Scr{-> Srch:} and \Scr{Replace with:};
with confirmation, \Scr{Replace? Yes (No)} at each occurrence.

\WPSide{Setup (Shift-F1)}
\begin{tabular}{@{}ll@{}}
3 Environment & Setup: Environment\\
\quad 1 Backup Options & 2 Original Document Backup, then \Typed{y} or \Typed{n}\\
\quad 8 Units of Measure & 1 Display and Entry of Numbers for Margins,\\
 & \quad Tabs, etc., or 2 Status Line Display;\\
 & then \Typed{i}, \Typed{c}, \Typed{p}, \Typed{w} or \Typed{u}\\
\end{tabular}
\par Choices printed in grey on the Setup screens do nothing.  Exit (F7) goes
back a screen; Cancel returns to the document.

\WPSide{Undelete (F1 on the Editing Screen)}
\Scr{Undelete 1 Restore; 2 Previous Deletion: 0}: 1 Restore, 2 Previous
Deletion (also Up); Down shows the later deletion.

\WPSide{File (the menu bar)}
Retrieve, Save, Exit (Ctrl-Q).  See \InRef{Pull-Down Menus}.

%% ---------------------------------------------------------------------------
\WPAppendix{H}{Limits}
\begin{tabular}{@{}lp{70mm}@{}}
\textbf{Item} & \textbf{Limit}\\
Document & a document file of up to 64K (65,536 bytes), its settings
included\\
Document name & 8 characters, a period and 3 characters (see
\InRef{Filenames})\\
Name prompt & 15 characters typed\\
Margin & less than 9 inches\\
Line spacing & less than 9 lines\\
Documents open & one\\
Search string & 256 bytes of the document file: a character, \Code{[HRt]}
or \Code{[HPg]} one each, other codes their length (Appendix D); the same
for the new text of Replace\\
Deletions kept & three, of up to 32K (32,768 bytes) each\\
\end{tabular}
